\documentclass[11pt]{article}
\usepackage[margin=1in]{geometry}
\usepackage[section]{placeins}
\usepackage{amsmath,amssymb,amsthm,mathtools}
\usepackage{booktabs,tabularx,array}
\usepackage{graphicx}
\usepackage{xcolor}
\usepackage{enumitem}
\usepackage[hidelinks]{hyperref}

\graphicspath{{../figures/}{figures/}}
\newcommand{\fillin}[1]{\textcolor{red}{\textbf{[#1]}}}
\makeatletter
\newcommand{\resdef}[2]{\expandafter\def\csname res@#1\endcsname{#2}}
\newcommand{\res}[1]{\ifcsname res@#1\endcsname\csname res@#1\endcsname\else\fillin{#1}\fi}
\makeatother
\IfFileExists{small_results.tex}{% src/small_example.py
% Generated automatically -- do not edit by hand.
\resdef{smThetaDiff}{5.6\times10^{-17}}
\resdef{smPostDiff}{2.2\times10^{-16}}
\resdef{smOddsDiff}{1.1\times10^{-16}}
\resdef{smPostSpam}{0.645073}
\resdef{smUnderflowAt}{162}
\resdef{smLogSum}{-2302.5851}
}{}
\IfFileExists{sim_results.tex}{% src/check_simulation.py
% Generated automatically -- do not edit by hand.
\resdef{simMatch}{100.0}
\resdef{simErrPrior}{5.6\times10^{-17}}
\resdef{simErrLik}{0}
\resdef{simErrJoint}{1.8\times10^{-15}}
\resdef{simErrPost}{1.8\times10^{-15}}
\resdef{simSmoothMaxZ}{4.73}
\resdef{simRareTheta}{8.4\times10^{-5}}
\resdef{simRareCount}{2}
\resdef{simRareBias}{2.6\times10^{-4}}
\resdef{simRareSE}{1.5\times10^{-4}}
\resdef{simTrueErrHam}{0.0078}
\resdef{simTrueErrSpam}{0.0216}
\resdef{simMaxZHam}{2.98}
\resdef{simMaxZSpam}{1.67}
\resdef{simWithin}{100}
\resdef{simNHam}{3789}
\resdef{simNSpam}{1011}
\resdef{simPiSpam}{0.2106}
\resdef{simPiZ}{1.06}
\resdef{simErrNFifty}{0.1201}
\resdef{simErrNFiveHundred}{0.0568}
\resdef{simErrNFiveThousand}{0.0162}
\resdef{simSlopeHam}{-0.48}
\resdef{simSlopeSpam}{-0.41}
\resdef{simLibMaxAllN}{0}
\resdef{simTN}{311}
\resdef{simFP}{15}
\resdef{simFN}{21}
\resdef{simTP}{53}
\resdef{simPrec}{0.779}
\resdef{simRec}{0.716}
\resdef{simPrecCI}{[0.667, 0.862]}
\resdef{simRecCI}{[0.605, 0.806]}
\resdef{simPairRatioHam}{1.002}
\resdef{simPairMedianHam}{1.000}
\resdef{simPairRatioSpam}{1.009}
\resdef{simPairMedianSpam}{0.997}
\resdef{simCalGap}{0.037}
\resdef{simExtreme}{0.117}
\resdef{simShareTr}{0.2106}
\resdef{simShareTe}{0.1850}
\resdef{simOOV}{0.0000}
\resdef{stressPrecZero}{0.837}
\resdef{stressRecZero}{0.764}
\resdef{stressRecHalf}{0.360}
\resdef{stressRecNine}{0.053}
}{}
\IfFileExists{real_results.tex}{% src/run_real.py
% Generated automatically -- do not edit by hand.
\resdef{realN}{5{,}574}
\resdef{realSpam}{747}
\resdef{realEmpty}{2}
\resdef{realTrain}{4{,}460}
\resdef{realTest}{1{,}114}
\resdef{realShareTr}{13.41}
\resdef{realShareTe}{13.38}
\resdef{realV}{7{,}818}
\resdef{realNHam}{55{,}091}
\resdef{realNSpam}{15{,}157}
\resdef{realMatch}{100.00}
\resdef{realErrPrior}{1.7\times10^{-16}}
\resdef{realErrLik}{0}
\resdef{realErrJoint}{3.6\times10^{-15}}
\resdef{realErrPost}{1.4\times10^{-14}}
\resdef{realRelJoint}{3.2\times10^{-18}}
\resdef{realSMin}{-1148.6}
\resdef{realTN}{960}
\resdef{realFP}{5}
\resdef{realFN}{14}
\resdef{realTP}{135}
\resdef{realPrec}{0.964}
\resdef{realRec}{0.906}
\resdef{realPrecCI}{[0.919, 0.985]}
\resdef{realRecCI}{[0.848, 0.943]}
\resdef{realAcc}{0.983}
\resdef{realBase}{0.866}
\resdef{realCMDiff}{0}
\resdef{realLeak}{133}
\resdef{realUnique}{5{,}171}
\resdef{dedupTest}{1{,}035}
\resdef{dedupTN}{899}
\resdef{dedupFP}{5}
\resdef{dedupFN}{13}
\resdef{dedupTP}{118}
\resdef{dedupPrec}{0.959}
\resdef{dedupRec}{0.901}
\resdef{dedupPrecCI}{[0.908, 0.983]}
\resdef{dedupRecCI}{[0.838, 0.941]}
\resdef{dedupAcc}{0.983}
\resdef{dedupMatch}{100.00}
\resdef{dedupErrJoint}{1.4\times10^{-14}}
\resdef{realB}{-1.865}
\resdef{realTopSpam}{\texttt{claim}, \texttt{prize}, \texttt{150p}, \texttt{won}, \texttt{cs}, \texttt{tone}}
\resdef{realTopHam}{\texttt{gt}, \texttt{lt}, \texttt{he}, \texttt{lor}, \texttt{she}, \texttt{da}}
\resdef{realTopSpamW}{5.60}
\resdef{realTopSpamWord}{\texttt{claim}}
\resdef{realPairMedianSpam}{0.96}
\resdef{realPairShareTwoSpam}{0.03}
\resdef{realPairCountSpam}{492}
\resdef{realPairMedianHam}{0.94}
\resdef{realPairShareTwoHam}{0.04}
\resdef{realPairCountHam}{2683}
\resdef{realPairRows}{spam & \texttt{free}, \texttt{call} & 2.18\times10^{-4} & 2.09\times10^{-4} & 1.05 \\
spam & \texttt{claim}, \texttt{prize} & 1.26\times10^{-4} & 3.16\times10^{-5} & 3.99 \\
spam & \texttt{txt}, \texttt{stop} & 9.51\times10^{-5} & 5.99\times10^{-5} & 1.59 \\
ham & \texttt{i}, \texttt{you} & 0.0011403 & 9.28\times10^{-4} & 1.23 \\
ham & \texttt{going}, \texttt{home} & 9.91\times10^{-6} & 6.09\times10^{-6} & 1.63 \\
ham & \texttt{i}, \texttt{love} & 2.12\times10^{-4} & 9.92\times10^{-5} & 2.14}
\resdef{realExtreme}{0.856}
\resdef{realWrongConf}{0.0063}
\resdef{realOOV}{0.0600}
\resdef{realLenCorr}{0.23}
\resdef{realRecShort}{0.815}
\resdef{realRecLong}{0.926}
\resdef{realNShort}{27}
\resdef{realNErr}{19}
}{}
\newcommand{\maybefig}[2]{%
  \IfFileExists{../figures/#1}{\includegraphics[width=#2]{#1}}{%
  \IfFileExists{figures/#1}{\includegraphics[width=#2]{#1}}{%
  \fbox{\parbox[c][4cm][c]{#2}{\centering\small Run the scripts\\
  to create \texttt{figures/\detokenize{#1}}}}}}}
\newcommand{\owner}[2]{\par\noindent\textit{Owner: #1 \hfill Checked by: #2}\par\medskip}

\theoremstyle{plain}
\newtheorem{theorem}{Theorem}
\newtheorem{proposition}[theorem]{Proposition}
\newcommand{\R}{\mathbb{R}}
\newcommand{\V}{\mathcal{V}}
\newcommand{\Vs}{|\mathcal{V}|}
\newcommand{\bx}{\mathbf{x}}
\newcommand{\bw}{\mathbf{w}}
\DeclareMathOperator*{\argmax}{arg\,max}
\DeclareMathOperator{\E}{E}
\DeclareMathOperator{\Cov}{Cov}

\title{DS5020 Final Project -- Iteration 3\\[2pt]
\large Naive Bayes Spam Classification: Small Worked Example, From-Scratch Implementation and Verification}
\author{Jeevan Mandala \and Anand Joel Battula}
\date{Fall 2026}

\setlength{\emergencystretch}{3em}
\begin{document}
\maketitle

\noindent\textbf{Project Leader (Iteration 3):} Anand Joel Battula\\
\textbf{Repository:} \url{https://github.com/Jeev2727/Ds5020-naive-bayes-classifier}\\
\textbf{Overleaf:} \href{https://www.overleaf.com/project/6b00d37c0b39a891f5669bfb/share#b636ef3bb6d1cd0b79b8563a9355bf45898555bc0a6e8363}{\texttt{Overleaf code}}

\medskip
\noindent Notation, definitions, assumptions A1--A7 and Theorems 1--10 are those of the Iteration 2
report; theorem numbers below refer to it unless they are new (Proposition~\ref{prop:pairs}).
Dataset, simulation model, parameter values and seed are unchanged from Iteration 2. The one
change is the tokenizer, which now deletes apostrophes as the Iteration 3 brief asks
(Section~\ref{sec:real}).

\tableofcontents

%======================================================================
\section{Team and Ownership}
Anand Joel Battula leads this iteration (he had not led before) and checked that each member
owns part of the math and part of the code.

\begin{center}
\begin{tabularx}{\textwidth}{@{}lXll@{}}
\toprule
Part & Content & Owner & Checker\\
\midrule
Math & Q1 worked example, Q2 underflow (\S\ref{sec:q1}--\ref{sec:q2}) & Jeevan & Anand\\
Math & Q3 library comparison, Q4 confusion matrix, Q5 conditional independence with Proposition~\ref{prop:pairs} (\S\ref{sec:q3}--\ref{sec:q5}) & Anand & Jeevan\\
Math & Simulation test and stress test analysis (\S\ref{sec:sim}) & Jeevan & Anand\\
Math & Assumption checks (\S\ref{sec:assump}) & Anand & Jeevan\\
Code & \texttt{methods.py}, \texttt{simulate.py}, \texttt{small\_example.py} & Jeevan & Anand\\
Code & \texttt{check\_simulation.py}, \texttt{run\_real.py}, \texttt{verify.py} & Anand & Jeevan\\
\bottomrule
\end{tabularx}
\end{center}

%======================================================================
\section{Code-to-Math Map}\label{sec:map}
Every numbered step of the Iteration 2 algorithm is one NumPy function in \texttt{src/methods.py},
with the equation in a comment above it. The same functions run the small example, the simulation
and the real data. scikit-learn and \texttt{numpy.linalg} appear only in the check scripts.

\begin{center}
\small
\begin{tabularx}{\textwidth}{@{}c>{\raggedright\arraybackslash}X l>{\raggedright\arraybackslash}p{4.3cm}@{}}
\toprule
Step & Equation & Source & Function in \texttt{methods.py}\\
\midrule
0 & tokens $w_1,\dots,w_L$; $\V$ from training; $X_{dw}=x_{dw}$ & Def.~9 & \texttt{my\_tokenize}\newline\texttt{my\_build\_vocab}\newline\texttt{my\_count\_matrix}\\
1 & $Y_{dc}=\mathbf 1\{y_d=c\}$;\ \ $C=Y^\top X$ & \S8 & \texttt{my\_one\_hot}\newline\texttt{my\_class\_counts}\\
2 & $\mathbf N=C\mathbf 1$ & \S8 & \texttt{my\_class\_totals}\\
3 & $\log\theta_{w\mid c}=\log(n_{w,c}+\alpha)-\log(N_c+\alpha\Vs)$ & Thm.~4 & \texttt{my\_log\_likelihoods}\\
4 & $\log\pi_c=\log\big((Y^\top\mathbf 1)_c/n\big)$ & Thm.~2 & \texttt{my\_log\_priors}\\
5 & $S=X(\log\Theta)^\top+\mathbf 1\log\boldsymbol\pi^\top$;\ \ $\bw=\log\boldsymbol\theta_1-\log\boldsymbol\theta_0$, $b=\log\frac{\pi_1}{\pi_0}$ & Thm.~6 & \texttt{my\_scores}\newline\texttt{my\_linear\_rule}\\
6 & $\hat y_d=\argmax_cS_{dc}$; posterior by log-sum-exp & Prop.~5, 9 & \texttt{my\_predict}\newline\texttt{my\_posterior}\\
7 & $CM=Y_{\text{true}}^\top Y_{\text{pred}}$; precision, recall; Wilson interval & Def.~10, Thm.~10 & \texttt{my\_confusion\_matrix}\newline\texttt{my\_precision\_recall}\newline\texttt{my\_wilson}\\
Q5 & $\hat P(w_1,w_2\mid c)=(X_c^\top X_c)_{w_1w_2}\big/\sum_dL_d(L_d-1)$ & Prop.~\ref{prop:pairs} & \texttt{my\_pair\_ratio\_matrix}\\
\bottomrule
\end{tabularx}
\end{center}
The wrappers \texttt{my\_nb\_fit} (steps 1--4) and \texttt{my\_nb\_predict} (steps 5--6) have the
interface used in the brief's starter code. Note that the confusion matrix is itself a matrix
product: entry $(i,j)$ of $Y_{\text{true}}^\top Y_{\text{pred}}$ counts messages with $y=i$ and $\hat y=j$.

%======================================================================
\section{Iteration 3 -- Math Questions}\label{sec:questions}

%----------------------------------------------------------------------
\subsection{Q1. Small worked example, step by step}\label{sec:q1}
Computed exactly with \texttt{fractions.Fraction} and, separately, with the floating-point
functions of \texttt{methods.py} (\texttt{src/small\_example.py}).

\textbf{Data.} $\V=(\text{free},\text{win},\text{call},\text{meet},\text{home})$, $\Vs=5$, $\alpha=1$.
\begin{center}
\begin{tabular}{@{}llccccc@{}}
\toprule
Message & Class & free & win & call & meet & home\\
\midrule
$d_1$ ``free win call'' & spam & 1 & 1 & 1 & 0 & 0\\
$d_2$ ``win free free'' & spam & 2 & 1 & 0 & 0 & 0\\
$d_3$ ``call win'' & spam & 0 & 1 & 1 & 0 & 0\\
$d_4$ ``meet home'' & ham & 0 & 0 & 0 & 1 & 1\\
$d_5$ ``call meet home meet'' & ham & 0 & 0 & 1 & 2 & 1\\
\midrule
new: ``free call home'' & ? & 1 & 0 & 1 & 0 & 1\\
\bottomrule
\end{tabular}
\end{center}

\textbf{Step 1: class counts} $C=Y^\top X$ (Iteration 2, \S8). The rows of $Y$ are
$(0,1)$ for $d_1,d_2,d_3$ and $(1,0)$ for $d_4,d_5$, so row 0 of $C$ sums the ham rows of $X$ and row 1 the spam rows:
\[
C=\begin{pmatrix}0&0&1&3&2\\3&3&2&0&0\end{pmatrix}\quad(\text{ham},\ \text{spam}).
\]
\textbf{Step 2: totals} $\mathbf N=C\mathbf 1=(6,\,8)^\top$.

\textbf{Step 3: smoothed likelihoods} (Theorem 4), $\theta_{w\mid c}=(n_{w,c}+1)/(N_c+5)$:
\[
\boldsymbol\theta_{\text{ham}}=\Big(\tfrac1{11},\tfrac1{11},\tfrac2{11},\tfrac4{11},\tfrac3{11}\Big),\qquad
\boldsymbol\theta_{\text{spam}}=\Big(\tfrac4{13},\tfrac4{13},\tfrac3{13},\tfrac1{13},\tfrac1{13}\Big).
\]
\textbf{Step 4: priors} (Theorem 2): $\pi_{\text{ham}}=2/5$, $\pi_{\text{spam}}=3/5$.

\textbf{Step 5a: classify with probabilities} (Theorem 1), using only the words that occur:
\[
\text{ham: }\tfrac25\cdot\tfrac1{11}\cdot\tfrac2{11}\cdot\tfrac3{11}=\tfrac{12}{6655},\qquad
\text{spam: }\tfrac35\cdot\tfrac4{13}\cdot\tfrac3{13}\cdot\tfrac1{13}=\tfrac{36}{10985}.
\]
The normalizer is $P(\bx)=\frac{12}{6655}+\frac{36}{10985}=\frac{14856}{2924207}$, so
\[
P(\text{spam}\mid\bx)=\frac{36/10985}{14856/2924207}=\frac{3993}{6190}\approx\res{smPostSpam},
\qquad P(\text{ham}\mid\bx)=\frac{2197}{6190}.
\]
Prediction: spam.

\textbf{Step 5b: classify with logs} (Theorem 6): $s_c=\log\pi_c+\sum_wx_w\log\theta_{w\mid c}$ gives
$s_{\text{ham}}=\log\frac{12}{6655}=-6.318217$ and $s_{\text{spam}}=\log\frac{36}{10985}=-5.720767$.
The larger log-score is spam, the same class as in Step 5a (Proposition 5).

\textbf{Step 6: posterior by log-sum-exp} (Proposition 9): with $M=s_{\text{spam}}$,
$P(\text{spam}\mid\bx)=1/(1+e^{s_{\text{ham}}-M})=0.645073$, equal to $3993/6190$.

\textbf{Linear rule.} $\bw=\log\boldsymbol\theta_{\text{spam}}-\log\boldsymbol\theta_{\text{ham}}
=(1.2192,\,1.2192,\,0.2384,\,-1.5533,\,-1.2657)$ and $b=\log\frac{3/5}{2/5}=\log\frac32=0.4055$. Then
$\bw^\top\bx+b=1.2192+0.2384-1.2657+0.4055=0.5974=\log\frac{3993}{2197}$, the log of the spam-to-ham
odds, and it is positive, so the rule predicts spam.

\textbf{Second-way checks.}
\begin{enumerate}[leftmargin=*,itemsep=1pt]
\item Rows of $\Theta$ sum to $11/11$ and $13/13$ exactly (Theorem 4).
\item The two posteriors sum to exactly 1: $2197/6190+3993/6190=1$.
\item Substituting back into Bayes' theorem: $\pi_{\text{spam}}P(\bx\mid\text{spam})/P(\bx)=3993/6190$.
\item The probability route and the log route predict the same class.
\item Code vs exact fractions: max $|\Theta|$ difference $\res{smThetaDiff}$, $|P(\text{spam}\mid\bx)|$
difference $\res{smPostDiff}$, log-odds difference $\res{smOddsDiff}$ (rounding level).
\end{enumerate}
\owner{Jeevan Mandala}{Anand Joel Battula}

%----------------------------------------------------------------------
\subsection{Q2. The underflow problem}\label{sec:q2}
Multiplying 500 probabilities of 0.01 should give $10^{-1000}$. In double precision the smallest
positive number is about $5\times10^{-324}$, so the running product becomes subnormal
($10^{-322}$ after 161 factors) and is exactly \texttt{0.0} after \res{smUnderflowAt} factors; the final
product is \texttt{0.0}. With logs the same score is
\[
\sum_{j=1}^{500}\log 0.01=500\log 0.01=\res{smLogSum},
\]
a perfectly ordinary number (log of a product is a sum of logs). The failure matters for
classification: if one class has probabilities 0.010 and the other 0.012, both products are
\texttt{0.0} and cannot be compared, while the log-scores $-2302.59$ and $-2211.42$ differ by
$500\log1.2=91.16$, a decisive margin. Real SMS messages are short (median 11 and 27 tokens),
but the log form costs nothing and also turns the product into the matrix product of Step 5.
\owner{Jeevan Mandala}{Anand Joel Battula}

%----------------------------------------------------------------------
\subsection{Q3. Comparison with \texttt{MultinomialNB(alpha=1)}}\label{sec:q3}
Both models are fitted on the \emph{same} count matrix, so they compute the same formula from the
same numbers.
\begin{center}
\small
\begin{tabular}{@{}lll@{}}
\toprule
Quantity & Simulation ($n=2000$) & Real data\\
\midrule
matching predictions & \res{simMatch}\% & \res{realMatch}\%\\
$\max|\log\hat\pi_{\text{ours}}-\log\hat\pi_{\text{lib}}|$ & $\res{simErrPrior}$ & $\res{realErrPrior}$\\
$\max|\log\Theta_{\text{ours}}-\log\Theta_{\text{lib}}|$ & $\res{simErrLik}$ & $\res{realErrLik}$\\
$\max|S_{\text{ours}}-S_{\text{lib}}|$ (joint log-prob) & $\res{simErrJoint}$ & $\res{realErrJoint}$\\
$\max|P(c\mid\bx)_{\text{ours}}-P(c\mid\bx)_{\text{lib}}|$ & $\res{simErrPost}$ & $\res{realErrPost}$\\
\bottomrule
\end{tabular}
\end{center}
\textbf{Why the differences are this size.} scikit-learn computes
$\log(C+\alpha)-\log(N+\alpha\Vs)$ with the same floating-point operations as Step 3, so $\log\Theta$
agrees to the last bit or within one rounding unit. The joint log-probabilities are sums of many
terms; adding them in a different order can change the last bit, and one unit in the last place of a
double is about $2.2\times10^{-16}$ times the size of the number (machine epsilon). On the real data
the scores reach $|\res{realSMin}|$, yet the largest difference is only $\res{realErrJoint}$, a relative
error of $\res{realRelJoint}$, below machine epsilon: the two implementations differ by less than one
rounding unit. This is far below the pass criterion $10^{-10}$, and there is no sign or scaling
convention to fix: unlike eigenvectors, these quantities are uniquely defined. The posteriors differ by up to
$\res{realErrPost}$ because $P(c\mid\bx)=e^{s_c-\log Z}$ amplifies rounding in the exponent, but this
is still nine orders of magnitude below the pass criterion.
\owner{Anand Joel Battula}{Jeevan Mandala}

%----------------------------------------------------------------------
\subsection{Q4. Confusion matrix, precision and recall}\label{sec:q4}
With $CM=Y_{\text{true}}^\top Y_{\text{pred}}=\begin{psmallmatrix}\mathrm{TN}&\mathrm{FP}\\ \mathrm{FN}&\mathrm{TP}\end{psmallmatrix}$
(our matrix equals \texttt{sklearn.metrics.confusion\_matrix}, max difference \res{realCMDiff}):
\begin{center}
\small
\begin{tabular}{@{}lccccc@{}}
\toprule
Test set & TN & FP & FN & TP & accuracy\\
\midrule
Real, all messages (\res{realTest}) & \res{realTN} & \res{realFP} & \res{realFN} & \res{realTP} & \res{realAcc}\\
Real, de-duplicated (\res{dedupTest}) & \res{dedupTN} & \res{dedupFP} & \res{dedupFN} & \res{dedupTP} & \res{dedupAcc}\\
Simulation (400) & \res{simTN} & \res{simFP} & \res{simFN} & \res{simTP} & --\\
\bottomrule
\end{tabular}

\medskip
\begin{tabular}{@{}lcc@{}}
\toprule
Test set & precision [95\% Wilson] & recall [95\% Wilson]\\
\midrule
Real, all messages & \res{realPrec}\ \res{realPrecCI} & \res{realRec}\ \res{realRecCI}\\
Real, de-duplicated & \res{dedupPrec}\ \res{dedupPrecCI} & \res{dedupRec}\ \res{dedupRecCI}\\
Simulation & \res{simPrec}\ \res{simPrecCI} & \res{simRec}\ \res{simRecCI}\\
\bottomrule
\end{tabular}
\end{center}
Precision is the estimate of $P(\text{spam}\mid\text{predicted spam})$, computed from the
$\mathrm{TP}+\mathrm{FP}$ flagged messages; recall estimates $P(\text{predicted spam}\mid\text{spam})$
from the $\mathrm{TP}+\mathrm{FN}$ spam messages (Proposition 7). The always-ham baseline has accuracy
\res{realBase} on the real test set and recall 0, which is why accuracy alone is not informative here.

\textbf{What this means for the data.} Of the \res{realTP}${}+{}$\res{realFP} messages the classifier
flags, \res{realTP} are spam: a flagged message is spam with estimated probability \res{realPrec}
(95\% interval \res{realPrecCI}). Of the \res{realTP}${}+{}$\res{realFN} spam messages it catches
\res{realTP}, so a spam message is caught with probability \res{realRec} \res{realRecCI}. The errors are
asymmetric: only \res{realFP} of the \res{realTN}${}+{}$\res{realFP} ham messages are flagged, which
matters because a lost real message costs more than a missed spam. The recall interval is wider
because it rests on the 149 spam messages rather than the 140 flagged ones and its centre is further
from 1. After removing duplicates, precision and recall change by less than 0.01 and stay inside the
original intervals, so the \res{realLeak} test messages that also appear in training do not drive the
result (assumption A3).
\owner{Anand Joel Battula}{Jeevan Mandala}

%----------------------------------------------------------------------
\subsection{Q5. Is conditional independence true?}\label{sec:q5}
For a pair of different words we compare the probability that two different token positions of a
class-$c$ message hold $w_1$ and $w_2$ with the product $P(w_1\mid c)P(w_2\mid c)$. Proposition~\ref{prop:pairs}
gives an estimator whose ratio to the product is 1 when Assumption A1 holds.

\begin{proposition}\label{prop:pairs}
Under Assumption 1, for a class-$c$ message of length $L$ and words $w_1\neq w_2$,
$\E[x_{w_1}x_{w_2}\mid L,c]=L(L-1)\,\theta_{w_1\mid c}\theta_{w_2\mid c}$. Hence
\[
\hat P(w_1,w_2\mid c)=\frac{\sum_{d:\,y_d=c}x_{dw_1}x_{dw_2}}{\sum_{d:\,y_d=c}L_d(L_d-1)}
=\frac{(X_c^\top X_c)_{w_1w_2}}{\sum_dL_d(L_d-1)}
\]
estimates $\theta_{w_1\mid c}\theta_{w_2\mid c}$, and the ratio
$\hat P(w_1,w_2\mid c)/(\hat\theta_{w_1\mid c}\hat\theta_{w_2\mid c})$ is close to 1 for large samples.
\end{proposition}
\begin{proof}
Given $L$ and $c$, the counts are multinomial (Definition 4) with $\E[x_w]=L\theta_w$ and, for
$w_1\ne w_2$, $\Cov(x_{w_1},x_{w_2})=-L\theta_{w_1}\theta_{w_2}$ (write $x_w=\sum_j\mathbf 1\{W_j=w\}$ and use
bilinearity of covariance: terms with the same position give $-\theta_{w_1}\theta_{w_2}$, since one
position cannot hold both words; terms with different positions are 0 by independence). By the
definition of covariance,
$\E[x_{w_1}x_{w_2}]=\Cov+\E[x_{w_1}]\E[x_{w_2}]=-L\theta_{w_1}\theta_{w_2}+L^2\theta_{w_1}\theta_{w_2}=L(L-1)\theta_{w_1}\theta_{w_2}$.
Summing over the class-$c$ messages (linearity of expectation) shows that the numerator has
expectation $\theta_{w_1}\theta_{w_2}\sum_dL_d(L_d-1)$. The sum $\sum_d x_{dw_1}x_{dw_2}$ is entry
$(w_1,w_2)$ of $X_c^\top X_c$ by the definition of matrix multiplication. By the law of large
numbers both $\hat P(w_1,w_2\mid c)$ and $\hat\theta_{w_1}\hat\theta_{w_2}$ converge to
$\theta_{w_1}\theta_{w_2}$, so the ratio converges to 1.
\end{proof}
Using document-level presence (``both words appear'') instead would \emph{not} give 1 even when A1
holds, because a fixed length $L$ makes presence events dependent; the token-pair estimator avoids this.

\textbf{Simulation} (A1 true by construction, 200{,}000 messages): ratio for the two most frequent
ham words \res{simPairRatioHam}, spam words \res{simPairRatioSpam}; median over frequent pairs
\res{simPairMedianHam} (ham) and \res{simPairMedianSpam} (spam). This is what ``passing'' looks like.

\textbf{Real data} (training messages):
\begin{center}
\small
\begin{tabular}{@{}llccc@{}}
\toprule
Class & Words & $\hat P(w_1,w_2\mid c)$ & $\hat\theta_{w_1\mid c}\hat\theta_{w_2\mid c}$ & ratio\\
\midrule
\res{realPairRows}\\
\bottomrule
\end{tabular}
\end{center}
Over the frequent pairs (expected pair count above 20; the real data are much sparser than the
200{,}000 simulated messages), the median ratio is \res{realPairMedianSpam}
in spam (\res{realPairCountSpam} pairs, share above 2: \res{realPairShareTwoSpam}) and
\res{realPairMedianHam} in ham (\res{realPairCountHam} pairs, share above 2: \res{realPairShareTwoHam}),
compared with 1.00 in the simulation. So the \emph{typical} pair is close to independent, while a
tail of pairs is strongly dependent. Medians slightly below 1 are expected when a class is a mixture
of sub-topics (prize spam and dating spam; work and family messages): words from the same sub-topic
co-occur more often than the class-level product predicts (ratio above 1), and words from different
sub-topics less often (ratio below 1). The model's single $\boldsymbol\theta_c$ per class averages
over these sub-topics.

\textbf{Answer: conditional independence does not hold exactly.} Some pairs are close to
independent: \texttt{free} and \texttt{call} have ratio 1.05, so knowing that one token of a spam
message is ``free'' barely changes the chance that another is ``call''. Others are strongly
dependent: \texttt{claim} and \texttt{prize} co-occur about 4 times as often as independence predicts,
because they belong to the same fixed phrases (``claim your prize''); \texttt{txt}/\texttt{stop}
(1.59) and, in ham, \texttt{i}/\texttt{love} (2.14) follow the same pattern. Ratios above 1 mean the
classifier counts the evidence of a phrase once per word, so $|\bw^\top\bx+b|$ is too large and the
posteriors are too extreme (Section~\ref{sec:assump}, A1). The argmax is still usually right, which is
why precision and recall remain high.
\owner{Anand Joel Battula}{Jeevan Mandala}

%======================================================================
\section{Simulation Test: Recovering the Known Truth}\label{sec:sim}
\textbf{Model.} Unchanged from Iteration 2: $\pi=(0.8,0.2)$, $\boldsymbol\theta_0,\boldsymbol\theta_1\sim
\mathrm{Dir}(\mathbf 1_{20})$, $L=3$, seed 0. \texttt{simulate\_nb} in \texttt{src/simulate.py} draws
the random numbers in the Iteration 2 order, so it reproduces the Iteration 2 data exactly
(e.g.\ $N_{\text{ham}}=\res{simNHam}$, $N_{\text{spam}}=\res{simNSpam}$, $\hat\pi_{\text{spam}}=\res{simPiSpam}$).

\textbf{Against the library.} \res{simMatch}\% of predictions match and all quantities agree to
$\res{simErrJoint}$ or better (Section~\ref{sec:q3}). At every sample size in Figure~\ref{fig:errn}
the largest difference from scikit-learn is $\res{simLibMaxAllN}$.

\textbf{Against the truth.} The maximum-likelihood estimate $\hat\theta_{w\mid c}=n_{w,c}/N_c$ has
standard error $\sqrt{\theta_{w\mid c}(1-\theta_{w\mid c})/N_c}$ (binomial proportion). With
$N_{\text{ham}}=\res{simNHam}$ and $N_{\text{spam}}=\res{simNSpam}$ training tokens, \res{simWithin}\% of
the 40 MLE estimates lie within 3 standard errors (largest $z$: \res{simMaxZHam} ham,
\res{simMaxZSpam} spam), and $\hat\pi_{\text{spam}}$ is \res{simPiZ} standard errors from 0.2. The
smoothed estimates are within $\|\hat{\boldsymbol\theta}_c-\boldsymbol\theta_c\|_2=\res{simTrueErrHam}$ (ham)
and \res{simTrueErrSpam} (spam). The spam row is less accurate because spam is only 20\% of the
messages, so $N_{\text{spam}}$ is about a quarter of $N_{\text{ham}}$ and its standard errors are about
twice as large ($\sqrt{4}=2$).

\textbf{Smoothing bias.} One smoothed ham estimate is \res{simSmoothMaxZ} standard errors from the
truth. Its true probability is $\res{simRareTheta}$ and it was seen \res{simRareCount} times. Laplace
smoothing moves it by $(n+1)/(N+\Vs)-n/N=\res{simRareBias}$, which is larger than its standard error
$\res{simRareSE}$. This is the bias of the Dirichlet posterior mean (Theorem 8) for a very rare word,
not sampling error, and it vanishes as $N_c\to\infty$.

\begin{figure}[htbp]
\centering
\maybefig{sim3_error_vs_n.png}{0.48\textwidth}\hfill
\maybefig{sim3_stress_precision_recall.png}{0.48\textwidth}
\caption{Left: error against the truth for $n=50$ to $15{,}000$ simulated messages (fitted log--log slopes
\res{simSlopeHam} ham, \res{simSlopeSpam} spam; theory $-1/2$). Right: stress test.}
\label{fig:errn}
\end{figure}

\textbf{Error against $n$.} The error at $n=50$, $500$, $5000$ is \res{simErrNFifty},
\res{simErrNFiveHundred} and \res{simErrNFiveThousand} (both classes combined). The fitted slopes
\res{simSlopeHam} and \res{simSlopeSpam} are close to $-1/2$, the rate $1/\sqrt{N_c}$ of the standard error.

\textbf{Stress test: classes that become identical.} We set
$\boldsymbol\theta_{\text{spam}}\leftarrow(1-t)\boldsymbol\theta_{\text{spam}}+t\boldsymbol\theta_{\text{ham}}$,
which breaks the premise that word distributions separate the classes. Recall falls from
\res{stressRecZero} at $t=0$ to \res{stressRecHalf} at $t=0.5$ and \res{stressRecNine} at $t=0.9$
(Figure~\ref{fig:errn}, right). Precision stays near 0.8 until $t\approx0.9$: as $\bw\to\mathbf 0$ the
prior term $b=\log(0.2/0.8)<0$ dominates, so the classifier flags fewer messages, but the few it flags
still have strongly spam-like words. At $t=1$ the classes are identical, almost nothing is flagged, and
precision falls to the base rate 0.2, since a flagged message is then spam only by chance.
\owner{Jeevan Mandala}{Anand Joel Battula}

%======================================================================
\section{Real Data}\label{sec:real}
\textbf{Cleaning steps} (all in \texttt{src/run\_real.py}):
\begin{enumerate}[leftmargin=*,itemsep=1pt]
\item Load the UCI file: \res{realN} messages, \res{realSpam} spam. \textbf{Rows removed: 0.}
\item Lower-case; delete apostrophes (\texttt{don't}$\to$\texttt{dont}); replace every other character
outside \texttt{[a-z0-9]} by a space; tokens are the remaining words. This removes all punctuation as
the brief asks; in Iteration 2 apostrophes were kept inside words, so the vocabulary changes slightly.
\res{realEmpty} messages contain no tokens and are kept (they are classified by the prior alone).
\item Stratified 80/20 split with seed 0 (the Iteration 2 split): \res{realTrain} training and
\res{realTest} test messages, spam share \res{realShareTr}\% and \res{realShareTe}\%.
\item Vocabulary from the training messages only: $\Vs=\res{realV}$; $N_{\text{ham}}=\res{realNHam}$,
$N_{\text{spam}}=\res{realNSpam}$ tokens.
\item Second run: remove exact duplicate texts first (\res{realUnique} unique messages), then split
and fit in the same way.
\end{enumerate}

The tokenizer change reduces the training vocabulary from 7{,}905 (Iteration 2) to \res{realV}
words, because forms such as \texttt{don't} and \texttt{dont} now coincide.

\textbf{What the linear rule learned.} The largest weights $w_j=\log\theta_{j\mid\text{spam}}-\log\theta_{j\mid\text{ham}}$
(words seen at least 10 times) are \res{realTopSpam}; the most negative are \res{realTopHam}
(Figure~\ref{fig:weights}). For example, one occurrence of \res{realTopSpamWord} multiplies the
spam-to-ham odds by $e^{\res{realTopSpamW}}\approx 270$. The bias is $b=\log(\pi_1/\pi_0)=\res{realB}$,
so a message with no informative words is classified as ham. The spam words are those of prize and
premium-rate messages (\texttt{claim}, \texttt{prize}, \texttt{won}, \texttt{guaranteed}, \texttt{150p},
\texttt{www}); the ham words are personal and conversational (\texttt{he}, \texttt{she}, \texttt{later},
\texttt{lor}, \texttt{da}). The two strongest ham words, \texttt{gt} and \texttt{lt}, are an artifact
of the dataset: the corpus replaces numbers and names in personal messages by the HTML codes
\texttt{\&lt;\#\&gt;}, and removing punctuation leaves the tokens \texttt{lt} and \texttt{gt}. They
are therefore genuine evidence of a personal message in this corpus, but would not carry over to new data.

\begin{figure}[htbp]
\centering
\maybefig{real_top_weights.png}{0.85\textwidth}
\caption{Words with the most negative and most positive weights $w$ in the decision rule $\bw^\top\bx+b$.}
\label{fig:weights}
\end{figure}

%======================================================================
\section{Verification Table}\label{sec:verify}
Every number is produced by a script in \texttt{src/}.
\begin{center}
\small
\begin{tabularx}{\textwidth}{@{}>{\raggedright\arraybackslash}p{2.9cm}>{\raggedright\arraybackslash}X>{\raggedright\arraybackslash}p{2.6cm}>{\raggedright\arraybackslash}p{4.4cm}c@{}}
\toprule
Data & Quantity & Compared with & Difference & Pass?\\
\midrule
Small example & $\Theta$, $P(\text{spam}\mid\bx)$ & exact fractions & $\res{smThetaDiff}$, $\res{smPostDiff}$ & \checkmark\\
Small example & probability vs log route & same class & identical & \checkmark\\
Simulation ($n=2000$) & predictions & \texttt{MultinomialNB} & \res{simMatch}\% match & \checkmark\\
Simulation ($n=2000$) & joint log-prob $S$ & \texttt{MultinomialNB} & $\max|\cdot|=\res{simErrJoint}$ & \checkmark\\
Simulation ($n=2000$) & $\hat\theta$ (MLE) & true $\theta$ & \res{simWithin}\% within 3 SE & \checkmark\\
Simulation ($n=50\ldots15000$) & error slope & $-1/2$ & \res{simSlopeHam}, \res{simSlopeSpam} & \checkmark\\
Simulation (Q5) & pair ratio & 1 & median \res{simPairMedianHam}, \res{simPairMedianSpam} & \checkmark\\
Simulation (stress) & recall at $t=0.9$ & falls toward 0 & \res{stressRecNine} & \checkmark\\
Real data & predictions & \texttt{MultinomialNB} & \res{realMatch}\% match & \checkmark\\
Real data & $\log\Theta$ & \texttt{MultinomialNB} & $\max|\cdot|=\res{realErrLik}$ & \checkmark\\
Real data & joint log-prob $S$ & \texttt{MultinomialNB} & $\max|\cdot|=\res{realErrJoint}$, rel.\ $\res{realRelJoint}$ & \checkmark\\
Real data & confusion matrix & \texttt{sklearn.metrics} & \res{realCMDiff} & \checkmark\\
Real data (dedup) & joint log-prob $S$ & \texttt{MultinomialNB} & $\max|\cdot|=\res{dedupErrJoint}$ & \checkmark\\
\bottomrule
\end{tabularx}
\end{center}
\owner{Anand Joel Battula}{Jeevan Mandala}

%======================================================================
\section{Assumption Checks}\label{sec:assump}
Each check runs on the simulation first, where the assumption holds by construction, so that we know
what a pass looks like.
\begin{center}
\small
\begin{tabularx}{\textwidth}{@{}>{\raggedright\arraybackslash}p{2.6cm}>{\raggedright\arraybackslash}X>{\raggedright\arraybackslash}p{2.5cm}>{\raggedright\arraybackslash}p{3.6cm}@{}}
\toprule
Assumption & Check & Simulation & Real data\\
\midrule
A1 conditional independence & pair ratio (Q5); share of posteriors beyond 0.999 or 0.001; calibration (Fig.~\ref{fig:calib}) & ratio $\approx1$; extreme \res{simExtreme}; max bin gap \res{simCalGap} & median ratio \res{realPairMedianSpam} (spam), \res{realPairMedianHam} (ham); extreme \res{realExtreme}; confidently wrong \res{realWrongConf}\\
A2 bag of words & read the misclassified messages (\texttt{figures/real\_errors.txt}) & holds by construction & fails for chatty spam and promotional words used personally\\
A3 i.i.d.\ messages & test texts also in training; de-duplicated rerun & holds by construction & \res{realLeak} leaked; dedup precision \res{dedupPrec}, recall \res{dedupRec}\\
A4 same distribution & spam share train/test & \res{simShareTr} / \res{simShareTe} & \res{realShareTr}\% / \res{realShareTe}\%\\
A5 length uninformative & correlation of $L$ with $\bw^\top\bx+b$ (spam); recall by length & $L=3$ fixed & corr.\ \res{realLenCorr}; recall \res{realRecShort} ($L<20$, $n=\res{realNShort}$) vs \res{realRecLong}\\
A6 closed vocabulary & share of test tokens not in $\V$ & \res{simOOV} & \res{realOOV}\\
A7 correct labels & most confident errors & holds by construction & mostly holds; 2 fragments, 1 doubtful label\\
\bottomrule
\end{tabularx}
\end{center}
\textbf{Findings on the real data.}
\begin{itemize}[leftmargin=*,itemsep=2pt]
\item \textbf{A1 fails, as expected.} Besides the pair ratios of Q5, \res{realExtreme} of the real test
posteriors are beyond 0.999 or 0.001, against \res{simExtreme} in the simulation where A1 holds.
The model is overconfident: in the lowest bin it predicts $P(\text{spam})\approx0.002$ on average,
but 1.3\% of those messages are spam, and \res{realWrongConf} of all test messages are wrong with
confidence above 0.999. The probabilities should therefore be read as a ranking, not as exact
frequencies; the decisions are still good.
\item \textbf{A3 holds in practice.} \res{realLeak} test messages also occur in training, but the
de-duplicated run gives almost the same precision and recall (Section~\ref{sec:q4}).
\item \textbf{A4 holds} by the stratified split (\res{realShareTr}\% vs \res{realShareTe}\% spam).
\item \textbf{A5 fails mildly.} Length carries information the model ignores: recall is
\res{realRecShort} on short spam ($L<20$, $n=\res{realNShort}$) but \res{realRecLong} on longer spam.
Short spam has fewer words to add evidence to $\bw^\top\bx$, and the prior term $b<0$ then wins.
\item \textbf{A6:} \res{realOOV} of test tokens (6\%) are not in the training vocabulary and are
dropped. This loses some evidence (often new numbers or misspellings) but cannot produce $\log0$.
\item \textbf{A2 (bag of words) explains most errors.} Reading the \res{realNErr} misclassified
messages (\texttt{figures/real\_errors.txt}): of the \res{realFN} missed spam messages, about half are
written like personal messages (``Hi babe its Chloe, how r u?'', ``Can U get 2 phone NOW?''), so their
words are ham words; what marks them as spam is the structure, a chatty opening followed by a
premium-rate sign-off, which word counts cannot see. Of the \res{realFP} false alarms, most contain
words that are usually promotional but are used in a personal sense: ``free download'', ``contact us'',
``customer service'', ``missed calls''.
\item \textbf{A6 interacts with the errors.} Spam sign-offs are often written without spaces
(``\texttt{Calls\pounds1/min\allowbreak moremobs\allowbreak EMSPOBox\allowbreak 45PO139WA}''), so cleaning produces long one-off tokens that are not
in $\V$ and are dropped, removing exactly the evidence that would have flagged the message.
\item \textbf{A7 (labels) is mostly fine, with two kinds of doubt.} Two spam messages are fragments
of longer messages (``\dots and is not connected to the'', ``accordingly. I repeat, just text the word
ok\dots''), so the label is right but the fragment carries almost no spam evidence; one spam message
(``Do you realize that in about 40 years, we'll have thousands of old ladies running around with
tattoos?'') contains no spam content at all and may be a forwarded joke labelled as spam.
\item \textbf{A1 again:} five spam messages receive $P(\text{spam}\mid\bx)<10^{-4}$, the overconfidence
predicted by the failure of conditional independence.
\end{itemize}

In the simulation the A4 shares differ (\res{simShareTr} vs \res{simShareTe}) because the Iteration 2
simulation splits by position (first 1600 / last 400), not by stratification; that difference is
sampling noise from the 400 test messages. In the real data the stratified split makes the shares
equal by design.

\begin{figure}[htbp]
\centering
\maybefig{sim3_reliability.png}{0.32\textwidth}\hfill
\maybefig{real_posterior_calibration.png}{0.64\textwidth}
\caption{A1 checks. Left: calibration on the simulation (points on the diagonal). Right: real data,
histogram of $P(\text{spam}\mid\bx)$ on the test set (log scale) and calibration.}
\label{fig:calib}
\end{figure}
\owner{Anand Joel Battula}{Jeevan Mandala}

%======================================================================
\section{Reproducibility}
From the repository root:
\begin{center}
\small
\begin{tabular}{@{}ll@{}}
\toprule
Command & Produces\\
\midrule
\texttt{python src/small\_example.py} & Q1, Q2; \texttt{docs/small\_results.tex}\\
\texttt{python src/check\_simulation.py} & simulation test; \texttt{sim3\_*.png}, \texttt{docs/sim\_results.tex}\\
\texttt{python src/run\_real.py} & real data; \texttt{real\_*.png}, \texttt{docs/real\_results.tex}\\
\bottomrule
\end{tabular}
\end{center}
The three \texttt{*\_results.tex} files hold every number in this report; uploading them to Overleaf
together with \texttt{figures/} updates the report without retyping anything.

\end{document}